% Generated draft: requires mathematical and source review.
\documentclass[11pt]{article}
\usepackage[margin=25mm]{geometry}
\usepackage{fontspec}
\setmainfont{Latin Modern Roman}
\usepackage{amsmath,amssymb,mathtools,unicode-math}
\setmathfont{Latin Modern Math}
\usepackage{enumitem,needspace,xurl,hyperref}
\setlength{\parindent}{0pt}
\setlength{\parskip}{5pt}
\setlength{\emergencystretch}{3em}
\newcommand{\N}{\mathbb{N}}
\newcommand{\Z}{\mathbb{Z}}
\newcommand{\Q}{\mathbb{Q}}
\newcommand{\R}{\mathbb{R}}
\newcommand{\C}{\mathbb{C}}
\newcommand{\F}{\mathbb{F}}
\newcommand{\A}{\mathbb{A}}
\newcommand{\PP}{\mathbb{P}}
\newcommand{\E}{\mathbb{E}}
\newcommand{\eps}{\varepsilon}
\newcommand{\dd}{\,\mathrm{d}}
\begin{document}
\section*{UnsolvedMath ID: 30004121}
\textbf{Source problem code:} catalogue entry
\subsection{Definitions and mathematical statement}
Use the standard extremal-graph convention that graphs are finite, simple, and undirected, and that containment means containment as a subgraph, not necessarily an induced subgraph. These conventions are implicit in the supplied notation and require confirmation against the source. For a fixed graph \(H\) with at least one edge and a positive integer \(n\), define
\[
\operatorname{ex}(n,H)=\max\bigl\{|E(G)|: |V(G)|=n\text{ and }G\text{ contains no subgraph isomorphic to }H\bigr\}.
\]
The rational exponents conjecture asks whether
\[
\begin{gathered}
\forall r\in\Q\text{ with }1<r<2,\quad
\exists\text{ a finite simple graph }H\text{ with }|E(H)|\geq 1,\\
\exists c,C\in\R\text{ with }0<c\leq C,\quad
\exists n_0\in\N,\ n_0\geq 1,\\
\forall n\in\N\text{ with }n\geq n_0,\qquad
c n^r\leq\operatorname{ex}(n,H)\leq C n^r.
\end{gathered}
\]
Here \(H\) may depend on \(r\), but is fixed independently of \(n\); the constants \(c,C,n_0\) may depend on \(r\) and \(H\). The target is one forbidden graph for each rational exponent, rather than a family of forbidden graphs. The endpoints \(r=1\) and \(r=2\) are excluded. The requirement that \(H\) have an edge makes the displayed maximum well-defined for every positive \(n\); an edgeless forbidden graph cannot realize the stated positive asymptotic growth.
\subsection{Short English statement}
For every rational exponent strictly between \(1\) and \(2\), is there one fixed forbidden graph \(H\) such that the maximum number of edges in an \(n\)-vertex graph containing no copy of \(H\) grows as \(\Theta(n^r)\)? The supplied record labels this conjecture open; this draft does not independently verify that status.
\subsection{Sources}
\begin{itemize}
\item Supplied local UnsolvedMath registry record, permanent ID \(30004121\), title ``Rational Exponents for Extremal Graph Functions,'' in the Oberwolfach Reports Open Problems collection; exact source problem code \texttt{catalogue entry}. The record lists proposed year \(2019\) and supplies no individual proposer.
\item Supplied report citation: \emph{Combinatorics, Probability and Computing} (2020), Oberwolfach Reports, pp.~1113--1189, DOI \texttt{10.4171/owr/2019/19}. \url{https://doi.org/10.4171/owr/2019/19}. The supplied background describes Oberwolfach Reports as documenting mathematical workshops and their open-problem discussions.
\item Kang--Kim--Liu, \emph{On the rational Turan exponents conjecture}. Supplied literature link: \url{https://doi.org/10.1016/j.jctb.2020.12.003}.
\item Jiang--Longbrake--Yepremyan, \emph{Rational exponents near 3/2}. Supplied literature link: \url{https://arxiv.org/abs/2607.19607}.
\item Literature assessment dated exactly \texttt{2026-08-22}: the supplied curation labels the single-graph rational exponents conjecture open and states that later work realizes broad infinite families of rational exponents without establishing realization by one forbidden graph for every rational in the interval. This is supplied dated literature triage derived from the OWR curation, not an independently verified assessment; its time-sensitive status and cited claims require independent verification.
\item The registry labels the statement \texttt{corrected\_verified} and says that notation, typography, and sentence structure were repaired while mathematical content was preserved. This is a supplied editorial verification claim, not mathematical review of this draft.
\end{itemize}
\end{document}
